% Caderno 4 — geometria da composição central rotacionável de dois fatores.
% Conceitual. O que a figura precisa ensinar é que, com alfa = raiz de 2,
% os oito pontos não centrais caem sobre uma mesma circunferência: é isso,
% e só isso, que significa "rotacionável".
% Contagem: N_E = 2^2 + 2*2 + 3 = 11 ensaios.
\documentclass[border=3pt,11pt]{standalone}
% Preâmbulo comum das figuras TikZ das apostilas.
% Carregado por todos os `standalone` deste diretório, para que os cinco
% diagramas não divirjam em fonte, cor ou espessura de traço.
%
% fontenc T1 + inputenc utf8 são obrigatórios: sem eles "ç", "ã" e "é"
% saem trocados ou desaparecem silenciosamente do PDF.
\usepackage[T1]{fontenc}
\usepackage[utf8]{inputenc}
\usepackage[brazil]{babel}
% A ordem importa: `amssymb` antes de `newtxmath`, exatamente como em
% `doeapostila.cls`. Invertida, as duas definem `\Bbbk` e a compilação aborta.
\usepackage{amsmath,amssymb}
\usepackage{newtxtext,newtxmath}   % mesma família serifada da classe doeapostila
\usepackage{tikz}
\usetikzlibrary{arrows.meta,calc,decorations.pathreplacing,positioning,
  shapes.geometric,backgrounds}

% Paleta: preto, cinzas e uma única cor de apoio.
\definecolor{azulDoE}{HTML}{1B4E79}
\definecolor{cinzaDoE}{gray}{0.55}
\definecolor{cinzaClaro}{gray}{0.82}
\definecolor{cinzaFundo}{gray}{0.93}

\tikzset{
  % vértice de plano fatorial, com a resposta escrita dentro
  vertice/.style={circle, draw=black, fill=white, line width=0.5pt,
                  inner sep=0pt, minimum size=8.5mm, font=\scriptsize},
  vertdestaque/.style={circle, draw=azulDoE, fill=azulDoE!10, text=azulDoE,
                  line width=0.9pt, inner sep=0pt, minimum size=8.5mm,
                  font=\scriptsize},
  % ponto de planejamento, sem texto dentro
  ponto/.style={circle, draw=black, fill=black, inner sep=0pt,
                minimum size=2.4mm},
  pontovazio/.style={circle, draw=black, fill=white, line width=0.5pt,
                inner sep=0pt, minimum size=2.6mm},
  pontoazul/.style={circle, draw=azulDoE, fill=azulDoE, inner sep=0pt,
                minimum size=2.8mm},
  axial/.style={rectangle, draw=azulDoE, fill=azulDoE, inner sep=0pt,
                minimum size=2.6mm},
  apagado/.style={circle, draw=cinzaClaro, fill=white, line width=0.5pt,
                inner sep=0pt, minimum size=2.6mm},
  aresta/.style={draw=black, line width=0.5pt},
  arestafundo/.style={draw=cinzaClaro, line width=0.5pt},
  eixo/.style={draw=black, line width=0.5pt, -{Latex[length=1.8mm,width=1.4mm]}},
  seta/.style={draw=black, line width=0.6pt, -{Latex[length=2.2mm,width=1.7mm]}},
  setaazul/.style={draw=azulDoE, line width=0.8pt,
                   -{Latex[length=2.2mm,width=1.7mm]}},
  cota/.style={draw=cinzaDoE, line width=0.5pt,
               {Latex[length=1.6mm]}-{Latex[length=1.6mm]}},
  % Corpos escolhidos para casar com as figuras de matplotlib, que usam
  % 8 pt nos ticks e 9 pt nos rótulos de eixo. Com a base de 11 pt do
  % `standalone` (a mesma da classe doeapostila), \scriptsize = 8 pt e
  % \footnotesize = 9 pt. Abaixo disso a figura deixa de ser legível na
  % página impressa, e é para isso que a folha de prova serve.
  rotulo/.style={font=\footnotesize},
  rotulinho/.style={font=\scriptsize},
  nota/.style={font=\footnotesize, align=left},
  notaazul/.style={font=\footnotesize, align=left, text=azulDoE},
}

\begin{document}
\begin{tikzpicture}[x=2.05cm, y=2.05cm]

  \def\alfa{1.41421}

  % --- circunferência de rotabilidade -------------------------------------
  \draw[cinzaDoE, dashed, line width=0.5pt] (0,0) circle (\alfa);

  % --- eixos ---------------------------------------------------------------
  \draw[eixo] (-1.95,0) -- (1.95,0) node[rotulo, right, yshift=-0.5mm] {$z_1$};
  \draw[eixo] (0,-1.95) -- (0,1.95) node[rotulo, above] {$z_2$};

  % --- quadrado fatorial ---------------------------------------------------
  \draw[cinzaClaro, line width=0.5pt]
       (-1,-1) -- (1,-1) -- (1,1) -- (-1,1) -- cycle;

  % --- raio até um ponto fatorial e até um ponto axial ---------------------
  \draw[cinzaDoE, line width=0.5pt] (0,0) -- (1,1);
  \node[rotulinho, text=cinzaDoE, anchor=south east, rotate=45]
       at (0.50,0.62) {$\sqrt{2}$};
  \draw[azulDoE, line width=0.7pt] (0,0) -- (\alfa,0);
  \node[rotulinho, text=azulDoE, anchor=north] at (0.72,-0.06)
       {$\alpha = \sqrt{2} \approx 1{,}414$};

  % --- pontos --------------------------------------------------------------
  \node[ponto] at (-1,-1) {};
  \node[ponto] at ( 1,-1) {};
  \node[ponto] at (-1, 1) {};
  \node[ponto] at ( 1, 1) {};

  \node[axial] at (-\alfa,0) {};
  \node[axial] at ( \alfa,0) {};
  \node[axial] at (0,-\alfa) {};
  \node[axial] at (0, \alfa) {};

  % três réplicas no centro, separadas o bastante para se contarem
  \node[pontoazul] at (-0.075,0.048) {};
  \node[pontoazul] at ( 0.075,0.048) {};
  \node[pontoazul] at ( 0,-0.082) {};

  % --- chamadas ------------------------------------------------------------
  \node[rotulinho, anchor=south west] at (1.05,1.04)
       {4 pontos fatoriais $(\pm 1, \pm 1)$};
  \node[rotulinho, text=azulDoE, anchor=south] at (0,\alfa+0.06)
       {4 pontos axiais, a $\pm\alpha$ sobre os eixos};
  \node[rotulinho, text=azulDoE, anchor=north east] at (-0.30,-0.34)
       {3 réplicas no centro};
  \draw[setaazul] (-0.27,-0.30) -- (-0.10,-0.11);

  \node[nota, anchor=north west] at (-1.95,-2.12)
    {$N_E = 2^{N_X} + 2N_X + N_{R0} = 4 + 4 + 3 = 11$ ensaios};
  \node[nota, anchor=north west] at (-1.95,-2.44)
    {Os oito pontos não centrais estão à mesma distância $\sqrt{2}$ da origem.};
  \node[notaazul, anchor=north west] at (-1.95,-2.76)
    {A variância da previsão passa a depender só da distância ao centro,
     não da direção.};

\end{tikzpicture}
\end{document}
